
\chapter{BACKGROUND}%This should be in all capital letters, "all caps!"

\section{Using this template}

To create your chapter 2, find the file named ``02Chapter.tex'' and open it in a latex editor. 
Then, replace all of this text. 
This file cannot be compiled directly; compiling this text and figures into a pdf is achevied by compiling the `OUThesis.tex' file.


\section{Figures}

Figures are an important part of any dissertation. 
\autoref{figureExample2} is an example figure. 
The caption in the code contains two parts, coded as ``\textbackslash caption[]\{\}''. 
Text placed between the square brackets [] goes to the `list of figures' in the front matter. 
Text placed between the curly brackets \{\} appears on the page with the figure. 

%\begin{figure}[tbp]
%\centering
%\includegraphics{figures/VectorBitmapExample}
%\caption[Title of figure that appears in the table of contents.]{Title of figure that appears in the text. 
%This figure is from Wikipedia, and shows the difference between a vector image and a bitmap. Note that if you zoom in all the way, the resolution is still perfect, even though the file size is just 46~kB}
%\label{figureExample2}
%\end{figure}

To get the figure to work properly, it should be the right size. 
Firstly, the width of the figure should be 6 inches or less. 
This will keep it out of the margins. 
The height of the figure + the height of the caption altogether must be small enough to fit within the margins.
This should be monitored manually. 

It is ideal to generate figures as vector images in pdf format. 
This will ensure infinite resolution with a small file size. 
If you are unsure what a vector image is, please research this on the internet. 
Also, when generating figures, make sure that all text in the figures uses a 12 point times font. 

OU formatting guidelines dictate that figures should be placed within a page of its first mention. 
How can this be achieved with LaTex?
Firstly, try putting the code for the figure just after it is mentioned. 
This is done with \autoref{figureExample1}.
Generally, LaTex is able to put the figure in a good location automatically.
Compile the file and see if its location is good. 

%\begin{figure}[tbp]
%\centering
%\includegraphics{figures/OregonRoute224}
%\caption[Oregon Route 224]{Oregon Route 224, vector image is 6 inches wide, and copied from wikipedia.}
%\label{figureExample1}
%\end{figure}

Very rarely, the location is not good.
In this case, the location can be forced. 
As the document will change continually, only make these changes immediately before the first formatting review. 
To force the location, the text in the code [tbp] should be changed. 
If it is changed to [b] that forces the figure to the bottom of the page. 
If it is changed to [p] that forces the figure to have its own page. 




\section{Code Figures}

There are times when code, specifically programming code, will be included in this thesis. 
Perhaps in the methods chapter? 
Perhaps elsewhere? 
As an example, see \autoref{code}.
\begin{figure}[bp]
\begin{lstlisting}[language=Python]
# This prints Hello, world!
print('Hello, world!')
print('Hello, friends!')
\end{lstlisting}
\caption[Coding figure; Hello world in python]{Coding figure; Hello world in python}
\label{code}
\end{figure}
This may require some changes or trouble shooting.

\section{Tables}

In latex, tables are efficient but also sort of a pain. 
\autoref{TableExample} shows an example of a table. 
There a many ways to format tables; please find internet resources. 
The only issue is that care must be taken to make sure your table is less than 6 inches wide. 
One tool that can help with this is ``\textbackslash parbox[]\{\}''.

\begin{table}[tbp]
\centering
\caption{Performance metrics for spectrum analysis}
\begin{tabular}{l c c l }
\toprule
{\bfseries Performance Metric} & {\bfseries Symbol} & {\bfseries Base Units} & {\bfseries Brief Description}\\
\midrule
Frequency accuracy & ~ & Hz & \parbox[t]{54 mm}{Accuracy with which the input\\signal frequency is identified.}\\
~ & ~ & ~ & \\
Scanning bandwidth & $\delta f_c$ & Hz & \parbox[t]{54 mm}{Total bandwidth to be\\analyzed.}\\
~ & ~ & ~ & \\
Scan rate & $\rho$ & Hz/s & \parbox[t]{54 mm}{The rate at which spectrum\\analysis is performed.}\\
\bottomrule
\label{TableExample}
\end{tabular}
\end{table}




